\section{On the Intersection Between Computational and Experimental Biology}

Molecular biology began in earnest in the early \num{20}th century with the discovery of proteins and DNA as the genetic material. In this short \num{100} years, the field has exploded. Compare this with many fields of mathematics, which date at least as far back as the 1700s. The first (intentional) sequence-specific DNA cutting was done by Hamilton Smith in 1978, and the Polymerase Chain Reaction was invented in 1983 by Kary Mullis. Entire genomes were first sequenced just before the turn of the century; the human genome was officially declared sequenced in 2003, after my birth. CRISPR-Cas9, in the form we know it today as a gene editing tool, was invented in 2012.

Clearly, it is no exaggerations to say that molecular biology is one of the fastest developing sciences today. Many things we take for granted are far younger than we may initially realize.

As we continue to move into the era of big data, the connections between biology and computer science will only continue to blur. Today, computers are used in genome sequencing, evolutionary biology, and, what I think has the biggest growth potential and is most relevant for this thesis, are molecular dynamics simulations. The ability to extrapolate from a protein structure its dynamics over time will have profound effects on the field of molecular biology in the decades to come. From protein-protein dynamics to drug design, biomolecular condensates, enzymatic activity and more, molecular dynamics simulations will be remembered as bringing upon yet another revolution in the field of molecular biology.

This is not to say that all of biology will become simulations, and experimental labs will close down in favor of computational ones. Rather, the two will co-exist with one another. A computational may suggest drugs predicted to have high binding affinity to a binding pocket, then the experimental labs must test them and report back for further compound optimization. A computational study might predict certain protein residues are essential in catalytic function of an enzyme, then a computational lab will test and verify this claim.

Working together, far richer models will be produced than ever could have been working independently.

\section{Thesis Aims}

It is the aim of this thesis to engage in the exact collaborative nature between computational and experimental biology discussed above.

Our protein of interest is p53, a crucial regulator of cellular signaling in times of stress. Previous computational work in our lab has made a number of predictions regarding the dynamics of p53 in mutations or to the binding of a small molecule. The work presented in the first half of this thesis extends these predictions to a wider variety of mutations, DNA sequences, and small molecules. Without a way to test them, however, they remain only predictions. That is why in the second half of this thesis, we will discuss an experimental assay in which we will be able to test and quantify the dynamics of p53 in the presence of these mutations, DNA sequences, and small molecules. Having set the stage for the first step in the computational to experimental pipeline, those who follow my work will be able to expand upon it and continue the computational-experimental-computational pipeline.

In this thesis, we set out with the following aims:

\begin{enumerate}
\item To study the effects of the Y220C mutant p53, and the restorative effects of PK11000 on the \num{12} p53 isoforms (Chapter \ref{chap:isoforms}).
\item To study the changing dynamics of p53 when bound to different DNA sequences important in a variety of cell-signaling pathways (Chapter \ref{chap:res}).
\item To study the effects of the \num{10} most common mutations to p53 associated with tumors (Chapter \ref{chap:mutants}).
\item To study the effects of the restorative compound Rezatapopt on the Y220C FL p53 and p53 DBD (Chapter \ref{chap:rezatapopt}).
\item To design an experimental assay in which predictions from the above chapters can be tested and verified (Chapter \ref{chap:beta-gal}).
\end{enumerate}